\chapter{Программная реализация}\label{ch:7}

В данном разделе описана реализация проекта. Разработка ведётся на языке Python 3 в среде разработки Visual Studio Code для единой работы с репозиторием - все конфигурационные файлы, код и тесты располагаются в единой структуре проекта, что упрощает сопровождение. 
Компоненты системы упаковываются в Docker-контейнеры и оркестрируются с помощью Docker Compose (контейнеры: PostgreSQL, Kafka с Zookeeper, воркеры машинного обучения, бэкенд, сервис бота, nginx). Данный подход позволяет воспроизвести рабочую среду как локально на машине разработчика, так и на сервере развёртывания.

\section{Программная реализация пайплайна обучения моделей и инференса}

В рамках пайплайна анализа изображения для решения подзадач нахождения деталей и повреждений автомобиля на фотографии используется семейство моделей YOLOv8 в режиме segment через библиотеку Ultralytics. Базовые веса задаются предобученной моделью yolov8m-seg.pt, затем обучение выполняется на прикладных данных.

\begin{figure}[h]
\centering
\includegraphics[scale=0.3]{details_cvat}
\caption{Пример разметки изображения для модели обнаружения деталей}
\end{figure}

Стек технологий для обучения моделей и инференса:

\begin{itemize}
    \item Python3 - реализация кода обучения и оценки
    \item Ultralytics (YOLOv8) - обучение и инференс моделей сегментации
    \item Numpy - быстрые численные операции в рамках приложения
    \item Pillow, OpenCV (opencv-python-headless) - загрузка изображений, предобработка
\end{itemize}

\begin{figure}[h]
\centering
\includegraphics[scale=0.5]{train_yolo}
\caption{Пример размеченных данных}
\end{figure}

Для обучения модели обнаружения и сегментации деталей было размечено порядка 250 изображений. 
На рисунке 18 представлен пример разметки деталей автомобиля на изображении. Экспорт размеченных данных производится в формате Ultralytics YOLO Segmentation 1.0. 
На рисунке 19 показано, как выглядят размеченные данные, поступающие на вход на обучение модели. 

Каждая строка размеченных данных соответствует одному объекту и содержит идентификатор класса и последовательность нормализованных координат (x, y), описывающих контур объекта (маску). Координаты заданы в относительной системе (в диапазоне [0, 1]) относительно ширины и высоты изображения.

Аналогичная разметка была проведена для обучения модели сегментации повреждений. Для данной модели было размечено ориентировочно 370 изображений. 
Пример размеченного обрезанного изображения детали представлен на рисунке 20.

\begin{figure}[h]
\centering
\includegraphics[scale=0.33]{damages_cvat}
\caption{Пример разметки изображения для модели обнаружения повреждений}
\end{figure}


\section{Программная реализация серверной части}

Серверная реализация строится на многослойной архитектуре и ориентирована на асинхронное выполнение. 
Стек бэкенда:
\begin{itemize}
    \item FastAPI - фреймворк для REST API
    \item uvicorn - ASGI-сервер для запуска приложения
    \item asyncio - библиотека для асинхронного выполнения кода
    \item asyncpg - для асинхронного доступа к PostgreSQL
    \item aiokafka - асинхронный клиент Apache Kafka
    \item MinIO (S3 API) - хранение исходных и аннотированных изображений
    \item Pydantic - удобная валидация данных и конфигурация (pydantic-settings)
    \item Loguru - структурированное логирование в сервисах
\end{itemize}

Реализованы и используются механизмы обеспечения отказоустойчивости и согласованности данных:
\begin{itemize}
    \item Обмен данными между бэкендом и ML-воркерами производится через топики Kafka
    \item Паттерн Transactional Outbox: фиксируются события и доставка в брокер сообщений, снимает риск рассогласования вида "запись в БД прошла, а сообщение в Kafka не дошло"
    \item Реализованы таймауты и отслеживание статусов каждой заявки
\end{itemize}


\section{Реализация интерфейса в мессенджере}

\begin{figure}[h]
\centering
\includegraphics[scale=0.5]{interface_damages}
\caption{Интерфейс бота Вконтакте}
\end{figure}

Сценарий для пользователя реализуется в боте Вконтакте с использованием фреймворка vkbottle - асинхронный long poll, обработка сообщений и callback кнопок клавиатуры. Бот не обращается к ML напрямую, а вызывает HTTP API бэкенда и реагирует на уведомления от Kafka.

В ручном режиме ("Ручной ввод") пользователь последовательно выбирает деталь и тип повреждения, а в ML-режиме - отправляет фотографию и производится автомтаический анализ. В рамках реализации учтены ограничения платформы Вконтакте на размер и состав inline-клавиатур - длинное меню при выборе деталей разбито на несколько сообщений.



\section{Тестирование}

Функциональное и юнит-тестирование реализовано с помощью pytest и pytest-asyncio - проверяются единицы поведения доменных сервисов, сценарии бэкенда, клиент бота, генерация клавиатуры и пр. Для проверки работы кода, связанного с контроллерами (БД), используются интеграционные тесты с использованием библиотеки testcontainers, позволяющей поднять PostgreSQL в контейнере во время работы теста.

Реализована и внедрена непрерывная интеграция (CI/CD) через Github Actions - при каждом залитии новых изменений в ветку код проверяется на стиль и статический анализ (ruff, flake8 и mypy), собирается окружение и прогоняются все тесты проекта с порогом покрытия 70\%.

Также было проведено нагрузочное тестирование - написан скрипт на Python3, который отправляет более 100 запросов в секунду, содержащих изображения, для проверки отказоустойчивости системы. Фиксируются кол-во успешно обработанных сценариев, время отклика, а также другие детали работы системы, такие как, например, заполнение очереди в брокере сообщений.